% Section 1.3, from lectures/L01/README.md: what you should be able to do afterwards, the
% questions to test yourself, and the next lecture.

\section{Review}
\label{c1:sec:review}

\needspace{6\baselineskip}
\subsection*{What you should be able to do now}
After this chapter, you should be able to:
\begin{itemize}
  \item Name the four pieces of an embedded Linux system, and say for any given file which one it
    belongs to.
  \item Walk the boot chain from power-on to the first userspace process, and say what each step
    does that the one before it could not.
  \item Say what \code{-kernel} skips, and therefore what this course does not teach you.
  \item Explain what a toolchain triplet names, and why a sysroot is not the same thing as the
    host's \file{/usr/include}.
  \item Answer the licence question for a given component: what must be published, to whom, and
    when.
  \item Say why an application that only makes system calls is not a derivative work of the kernel,
    and why a kernel module usually is.
  \item Predict whether a module will load, given its \code{MODULE_LICENSE} and the symbols it uses.
\end{itemize}

\needspace{6\baselineskip}
\subsection*{Questions to test yourself}
\begin{enumerate}
  \item A colleague says ``we ship Linux, so all our code has to be open source''. What is wrong
    with that sentence, and what is the part of it that is right?
  \item Your product has a bootloader, a kernel, a BusyBox root filesystem and one proprietary
    daemon you wrote. A customer asks for the source. What do you have to give them?
  \item Why does the boot chain have an SPL at all, given that U-Boot could do the same job?
  \item What does \code{aarch64-linux-gnu} tell you that \code{aarch64} alone does not?
  \item A driver is licensed MIT and calls a function exported with \code{EXPORT_SYMBOL_GPL}. What
    happens, and at which point: compile, link, load, or run?
  \item Why is the kernel image 24 MB when the machine it boots has one serial port and one timer?
\end{enumerate}

\needspace{6\baselineskip}
\subsection*{Where to look}
\begin{itemize}
  \item The build targets are documented in the repository's root \file{README.md}, and
    \code{make help} lists them all.
\end{itemize}

\needspace{6\baselineskip}
\subsection*{Looking ahead}
Chapter~\ref{c2:ch} is about:
\begin{itemize}
  \item The command line, taken as the set of questions you are able to ask a running kernel.
  \item \file{/proc} and \file{/sys}: not directories, and not files either.
  \item What a two-megabyte userland gives up, measured against the one on your laptop.
  \item Why PID 1 is different from every other process, and what happens if it exits.
\end{itemize}
The machine you built in this chapter is the machine the other eleven chapters run on.
Chapter~\ref{c2:ch} is about learning to ask it what it is doing.
